% GENERATED FILE. Edit canonical lessons, then run npm run generate:books.
% canonical-source-hash: fnv1a64:ca04a32f895be43e
% canonical-lessons: MR-C91-jina, MR-C91-arsa, MR-C91-diva, MR-C91-pankha, MR-C91-killi

\chapter{Stairs, Mirror, Lamp, Fan, Key}
\label{ch:mr-stairs-mirror-lamp-fan-key}
\hlchaptermodality{91}

\begin{chapteropening}
\textbf{By the end of this chapter:} \emph{I can say stairs, a mirror, a lamp, a fan and a key.}

Complete the last lesson of chapter 91: i can say stairs, a mirror, a lamp, a fan and a key.
\end{chapteropening}

\section[jinā]{\mr{जिना} (jinā) — stairs}

\label{lesson:MR-C91-jina}



\begin{quote}
Before the new one: say the Marathi for a cupboard, then the Marathi for the ground.
\end{quote}

\subsection*{You'll want to know: \mr{जिना}}
\textbf{\mr{जिना}} — \emph{jinā} — \textquotedblleft{}stairs\textquotedblright{}.

\begin{grammarlens}[title={What you've built}]
One more word for this chapter.
\end{grammarlens}

\subsection*{Your turn}
\begin{itemize}
\raggedright
  \item \emph{Say it:} \emph{jinā}
  \item \emph{Say it:} \emph{jinā}, once more
  \item \emph{Say it:} say \emph{jinā}
  \item \emph{Recall:} say the Marathi for a cupboard, then the Marathi for the ground, then say \emph{jinā} again
\end{itemize}

\begin{tcolorbox}[breakable,colback=teal!4,colframe=teal!35!black,title={Before you move on}]
What does \mr{जिना} mean? (\textquotedblleft{}Stairs\textquotedblright{}.) Say it once more.
\end{tcolorbox}
\section[ārsā]{\mr{आरसा} (ārsā) — a mirror}

\label{lesson:MR-C91-arsa}



\begin{quote}
Before the new one: say the Marathi for the ground, then the Marathi for stairs.
\end{quote}

\subsection*{You'll want to know: \mr{आरसा}}
\textbf{\mr{आरसा}} — \emph{ārsā} — \textquotedblleft{}a mirror\textquotedblright{}.

\begin{grammarlens}[title={What you've built}]
One more word for this chapter.
\end{grammarlens}

\subsection*{Your turn}
\begin{itemize}
\raggedright
  \item \emph{Say it:} \emph{ārsā}
  \item \emph{Say it:} \emph{ārsā}, once more
  \item \emph{Say it:} say \emph{ārsā}
  \item \emph{Recall:} say the Marathi for the ground, then the Marathi for stairs, then say \emph{ārsā} again
\end{itemize}

\begin{tcolorbox}[breakable,colback=teal!4,colframe=teal!35!black,title={Before you move on}]
What does \mr{आरसा} mean? (\textquotedblleft{}A mirror\textquotedblright{}.) Say it once more.
\end{tcolorbox}
\section[divā]{\mr{दिवा} (divā) — a lamp}

\label{lesson:MR-C91-diva}



\begin{quote}
Before the new one: say the Marathi for stairs, then the Marathi for a mirror.
\end{quote}

\subsection*{You'll want to know: \mr{दिवा}}
\textbf{\mr{दिवा}} — \emph{divā} — \textquotedblleft{}a lamp\textquotedblright{}.

\begin{grammarlens}[title={What you've built}]
One more word for this chapter.
\end{grammarlens}

\subsection*{Your turn}
\begin{itemize}
\raggedright
  \item \emph{Say it:} \emph{divā}
  \item \emph{Say it:} \emph{divā}, once more
  \item \emph{Say it:} say \emph{divā}
  \item \emph{Recall:} say the Marathi for stairs, then the Marathi for a mirror, then say \emph{divā} again
\end{itemize}

\begin{tcolorbox}[breakable,colback=teal!4,colframe=teal!35!black,title={Before you move on}]
What does \mr{दिवा} mean? (\textquotedblleft{}A lamp\textquotedblright{}.) Say it once more.
\end{tcolorbox}
\section[paṅkhā]{\mr{पंखा} (paṅkhā) — a fan}

\label{lesson:MR-C91-pankha}



\begin{quote}
Before the new one: say the Marathi for a mirror, then the Marathi for a lamp.
\end{quote}

\subsection*{You'll want to know: \mr{पंखा}}
\textbf{\mr{पंखा}} — \emph{paṅkhā} — \textquotedblleft{}a fan\textquotedblright{}.

\begin{grammarlens}[title={What you've built}]
One more word for this chapter.
\end{grammarlens}

\subsection*{Your turn}
\begin{itemize}
\raggedright
  \item \emph{Say it:} \emph{paṅkhā}
  \item \emph{Say it:} \emph{paṅkhā}, once more
  \item \emph{Say it:} say \emph{paṅkhā}
  \item \emph{Recall:} say the Marathi for a mirror, then the Marathi for a lamp, then say \emph{paṅkhā} again
\end{itemize}

\begin{tcolorbox}[breakable,colback=teal!4,colframe=teal!35!black,title={Before you move on}]
What does \mr{पंखा} mean? (\textquotedblleft{}A fan\textquotedblright{}.) Say it once more.
\end{tcolorbox}
\section[killī]{\mr{किल्ली} (killī) — a key}

\label{lesson:MR-C91-killi}



\begin{quote}
Before the new one: say the Marathi for a lamp, then the Marathi for a fan.
\end{quote}

\subsection*{You'll want to know: \mr{किल्ली}}
\textbf{\mr{किल्ली}} — \emph{killī} — \textquotedblleft{}a key\textquotedblright{}.

\begin{grammarlens}[title={What you've built}]
One more word for this chapter.
\end{grammarlens}

\subsection*{Your turn}
\begin{itemize}
\raggedright
  \item \emph{Say it:} \emph{killī}
  \item \emph{Say it:} \emph{killī}, once more
  \item \emph{Say it:} say \emph{killī}
  \item \emph{Recall:} say the Marathi for a lamp, then the Marathi for a fan, then say \emph{killī} again
\end{itemize}

\begin{tcolorbox}[breakable,colback=teal!4,colframe=teal!35!black,title={Before you move on}]
What does \mr{किल्ली} mean? (\textquotedblleft{}A key\textquotedblright{}.) Say it once more.
\end{tcolorbox}
